% Options for packages loaded elsewhere
\PassOptionsToPackage{unicode}{hyperref}
\PassOptionsToPackage{hyphens}{url}
\documentclass[
  11pt,
]{article}
\usepackage{xcolor}
\usepackage[margin=1in]{geometry}
\usepackage{amsmath,amssymb}
\setcounter{secnumdepth}{-\maxdimen} % remove section numbering
\usepackage{iftex}
\ifPDFTeX
  \usepackage[T1]{fontenc}
  \usepackage[utf8]{inputenc}
  \usepackage{textcomp} % provide euro and other symbols
\else % if luatex or xetex
  \usepackage{unicode-math} % this also loads fontspec
  \defaultfontfeatures{Scale=MatchLowercase}
  \defaultfontfeatures[\rmfamily]{Ligatures=TeX,Scale=1}
\fi
\usepackage{lmodern}
\ifPDFTeX\else
  % xetex/luatex font selection
\fi
% Use upquote if available, for straight quotes in verbatim environments
\IfFileExists{upquote.sty}{\usepackage{upquote}}{}
\IfFileExists{microtype.sty}{% use microtype if available
  \usepackage[]{microtype}
  \UseMicrotypeSet[protrusion]{basicmath} % disable protrusion for tt fonts
}{}
\makeatletter
\@ifundefined{KOMAClassName}{% if non-KOMA class
  \IfFileExists{parskip.sty}{%
    \usepackage{parskip}
  }{% else
    \setlength{\parindent}{0pt}
    \setlength{\parskip}{6pt plus 2pt minus 1pt}}
}{% if KOMA class
  \KOMAoptions{parskip=half}}
\makeatother
\setlength{\emergencystretch}{3em} % prevent overfull lines
\providecommand{\tightlist}{%
  \setlength{\itemsep}{0pt}\setlength{\parskip}{0pt}}
\usepackage{bookmark}
\IfFileExists{xurl.sty}{\usepackage{xurl}}{} % add URL line breaks if available
\urlstyle{same}
\hypersetup{
  pdftitle={Characterizing the Feasible Region for Self-Modifying Substrates in Interactive Environments},
  pdfauthor={Avir Research},
  hidelinks,
  pdfcreator={LaTeX via pandoc}}

\title{Characterizing the Feasible Region for Self-Modifying Substrates
in Interactive Environments}
\author{Avir Research}
\date{2026-03-19}

\begin{document}
\maketitle

\emph{The shared artifact. Birth writes theory. Experiment writes
results. Compress edits both.}

\subsection{Abstract}\label{abstract}

We formalize six rules (R1-R6) for recursive self-improvement as
mathematical conditions on a state-update function
\(f: S \times X \to S\) and derive necessary properties of any system
satisfying all six simultaneously. From 505 experiments across 9
architecture families on ARC-AGI-3 interactive games, we extract 26
constraints and prove: (1) no satisfying system has a fixed point ---
self-modification is necessary, not optional; (2) in finite
environments, the system must process its own internal state to maintain
irredundant growth (the self-observation requirement); (3) the feasible
region is non-empty for Level 1 navigation but currently unoccupied for
the full constraint set including R3 (self-modification of operations).
Whether a substrate exists inside all six walls remains open. The
contribution is the walls themselves.

\subsection{1. Introduction}\label{introduction}

505 experiments across 9 architecture families (codebook/LVQ, reservoir,
graph, LSH, kd-tree, cellular automata, LLM, L2 k-means, Bloom filter)
tested substrates for navigation and classification in ARC-AGI-3
interactive games. All experiments used the same evaluation framework
(R1-R6) and constraint map (U1-U26).

The experiments carved a feasible region --- the set of systems that
satisfy all constraints simultaneously. This paper formalizes the
constraints mathematically, derives necessary properties of the feasible
region, and states honestly what is proven vs conjectured vs open.

\textbf{The paper is valid regardless of outcome.} If a substrate
satisfying all six rules is found, the feasible region is non-empty and
the characterization guided the search. If no such substrate exists, the
characterization identifies which constraints are mutually exclusive ---
also a contribution. We report whichever is true at time of writing.

\subsection{2. Related Work}\label{related-work}

\subsubsection{2.1 Self-referential
self-improvement}\label{self-referential-self-improvement}

Schmidhuber's Gödel Machine (2003) is the closest formal framework: a
self-referential system that rewrites its own code when it can prove the
rewrite is useful. Key differences from our framework: (1) Gödel
machines require a utility function (external objective --- our R1
prohibits this), (2) provable self-improvement is limited by Gödel's
incompleteness theorem, (3) the Gödel machine framework doesn't address
exploration saturation in finite environments.

\subsubsection{2.2 Intrinsic motivation and curiosity-driven
exploration}\label{intrinsic-motivation-and-curiosity-driven-exploration}

Pathak et al.~(2017) formalize curiosity as prediction error in a
learned feature space. Count-based exploration (Bellemare et al., 2016)
uses state visitation counts. Both address exploration in sparse-reward
environments. Key difference: these methods add intrinsic REWARD
signals, which function as objectives. Our R1 requires no objectives.
Our derivation shows self-observation is required by the constraint
system itself, not as a reward mechanism.

\subsubsection{2.3 Autopoiesis}\label{autopoiesis}

Maturana \& Varela (1972) define autopoietic systems as networks that
produce the components that produce the network. Organizationally
closed. Related to our fixed-point conjecture (Section 4.4). Key
difference: autopoiesis maintains structure (homeostasis); our U17
requires unbounded growth. Our system is autopoietic + growth.

\subsubsection{2.4 Growing neural gas and self-organizing
maps}\label{growing-neural-gas-and-self-organizing-maps}

Fritzke (1995) GNG, Kohonen (1988) SOM/LVQ. The codebook substrate is
LVQ + growing codebook. Well-characterized in the literature. Our
contribution is not the architecture --- it's the constraint map
extracted from systematic testing.

\subsection{3. Formal Framework}\label{formal-framework}

\subsubsection{3.1 Primitives}\label{primitives}

The substrate is a triple (f, g, F): - f\_s: X → S (state update
parameterized by current state s) - g: S → A (action selection) - F: S →
(X → S) (meta-rule mapping states to update rules; F IS the frozen
frame)

Dynamics: s\_\{t+1\} = F(s\_t)(x\_t), a\_t = g(s\_t).

\subsubsection{3.2 Structural Rules (R1-R6) as Conditions on
F}\label{structural-rules-r1-r6-as-conditions-on-f}

{[}From BIRTH.md Formalization 1 --- to be migrated and refined with
literature citations{]}

\subsubsection{3.3 Growth Topology (U3, U17, U20,
R6)}\label{growth-topology-u3-u17-u20-r6}

{[}From BIRTH.md Formalization 2 --- to be migrated and refined{]}

\subsection{4. Results}\label{results}

\subsubsection{4.1 The Core Tension (R3 + U7 + U17 +
U22)}\label{the-core-tension-r3-u7-u17-u22}

The system has no fixed point (Theorem 1). Self-modification is
necessary, not optional.

{[}From BIRTH.md Formalization 1, Section 1.2-1.3{]}

\subsubsection{4.2 Irredundant Growth}\label{irredundant-growth}

Every new component covers unique territory (Propositions 1-4).

{[}From BIRTH.md Formalization 2, Section 2.2{]}

\subsubsection{4.3 The Self-Observation
Requirement}\label{the-self-observation-requirement}

\textbf{Theorem 2:} In finite environments, U17 + R6 + R1 require
self-observation.

{[}From BIRTH.md Formalization 3 --- the central result{]}

\subsubsection{4.4 Fixed-Point Conjecture}\label{fixed-point-conjecture}

The minimal self-observing substrate is a fixed point of F. Conjectured,
not proven.

{[}From BIRTH.md Formalization 3, Section 3.4{]}

\subsection{5. Experimental Evidence}\label{experimental-evidence}

\subsubsection{5.1 Navigation (505
experiments)}\label{navigation-505-experiments}

All 3 ARC-AGI-3 games solved at Level 1: - LS20: LSH or k-means, 4
actions, argmin. 9/10 at 120K steps. - FT09: k-means, 69 actions (64
grid + 5 simple), argmin. 3/3. (Step 503) - VC33: k-means, 3 actions
(zone discovery), argmin. 3/3. (Step 505)

Unifying mechanism: graph + edge-count argmin + correct action
decomposition.

\subsubsection{5.2 Level 2 Failure as Feasibility
Violation}\label{level-2-failure-as-feasibility-violation}

259-cell plateau (Steps 486-492). Edge counts grow (U17 formally
satisfied) but each marginal count is functionally redundant (R6
violated). The system exits the feasible region after exploration
saturates.

\subsubsection{5.3 Architecture Family
Summary}\label{architecture-family-summary}

{[}9 families, kill reasons, experiment counts --- from
CONSTRAINTS.md{]}

\subsection{6. Degrees of Freedom}\label{degrees-of-freedom}

{[}Enumerated from birth derivation --- these become the next experiment
phase{]}

\subsection{7. Discussion}\label{discussion}

{[}To be written{]}

\subsection{References}\label{references}

\begin{itemize}
\tightlist
\item
  Schmidhuber, J. (2003). Gödel Machines: Self-Referential Universal
  Problem Solvers Making Provably Optimal Self-Improvements.
  arXiv:cs/0309048.
\item
  Pathak, D. et al.~(2017). Curiosity-driven Exploration by
  Self-Supervised Prediction. ICML.
\item
  Maturana, H. \& Varela, F. (1972). Autopoiesis and Cognition: The
  Realization of the Living.
\item
  Fritzke, B. (1995). A Growing Neural Gas Network Learns Topologies.
  NeurIPS.
\item
  Kohonen, T. (1988). Self-Organization and Associative Memory.
  Springer.
\item
  Bellemare, M. et al.~(2016). Unifying Count-Based Exploration and
  Intrinsic Motivation. NeurIPS.
\end{itemize}

\end{document}
